% 来源：OEBPS/Text/Section0001_split_001.xhtml#sigil_toc_id_14

\BookChapter{5}
\section[第5章\q 打破误区：粉碎9种投资营销谎言]{第5章\\ 打破误区：粉碎9种投资营销谎言}

\begin{BookQuote}
牢记黄金法则：谁有黄金，谁定法则。

\hfill ——无名氏
\end{BookQuote}

\begin{BookQuote}
你必须先学习游戏法则，然后你必须比其他人玩得更好。

\hfill ——阿尔伯特·爱因斯坦
\end{BookQuote}

我知道，你想直接跳过去学习投资，弄清楚自己的钱该被投到什么地方，以获得财务自由，我也想演示给你看。当我看到有人真的“懂了”，能够理解和接受投资理财这个游戏而且确实也可以赢，我会喜形于色。但是只是给你省钱，得到很高的回报，减少投资风险，这些并不够。你必须知道，这里有很多人想要拿走你的一部分财富。金融体系到处都是漏洞，我更想称之为地雷，它们会毁掉你的财务未来。所以在这个部分，我们将一起来击破9种投资营销谎言——证券公司、基金公司、银行等金融机构用这些谎话营销了很多年。如果你还没有意识到，或者如果你根本没有意识到它们把你洗脑了，这些投资谎言就会一步一步、有计划地摧毁你的财务未来。

从这部分开始，本书开始给你带来实实在在的好处了！如果你的工资处于美国平均水平，年收入5万美元，现在每年能够储蓄你工资收入的10\%，并用这些钱长期投资，那么你的整个投资生涯可以节省25万美元。你只要用到你在这个部分学到的一部分知识就可以做到了。你根本不用工作，就能得到你目前5年的收入，足够你按照你现在的生活方式过上5年。这一点是被实际统计数据证实的，并非我胡乱编造出来的。如果你的收入低于美国平均水平，每年只能挣到3万美元，每年你只能储蓄个人收入的5\%，那么你的整个投资生涯还是可以节省出来15万美元的，而你也不用工作就能拿到这么多钱。如果你的年收入高于美国平均水平，超过10万美元，那么通过学习这部分的投资理财知识，你的一生就能节省50万\textasciitilde{}100万美元。这听起来像是在吹牛，对不对？在本书后面的内容里，我会让数字来说话。

这一章很短，所以你一定要用心地读，因为你想马上采取行动。彻底粉碎这些投资理财谎言，你会立即“止血”。在原来你从来没有想过需要止血的地方，你可以省掉很多不必要的费用。了解了这9种投资营销谎言能保护你少受损失，确保你能够达到财务自由的水平，这正是我承诺要达到的理财目标。现在让我们开始吧。

\subsection{欢迎来到丛林}

不管你是一个经验老到的投资者，还是刚刚开始把自己看作投资者的投资菜鸟，面对的都是同样的金融世界，瑞·达利欧生动地把它描述为金融丛林，对于我们所有人来说，丛林中都布满了同样多的危险。但是这些危险大多数都隐藏在你所不知道的，会伤害你的事实里面。

\subsection{报价}

我希望你能这样想象一下，有人来找你谈生意，他说有下面这样的投资机会，想问问你有没有兴趣投资：他希望你投入100\%的资金，承担100\%的风险。如果这笔投资赚钱了，他要分走利润的60\%，甚至更多，作为他个人的管理费。顺便说一下，如果这笔投资亏钱的话，是你亏钱，你还得一样付他的投资管理费。

你干吗？

我敢肯定，这种事想都不用想。你的本能反应是：“我才不会做这种傻事呢！这也太傻了。”唯一的问题是，如果你像90\%的美国投资者一样，你投资的是典型的公募基金，不管你信还是不信，你其实已经同意签订上面那种投资合同条款了。

就是按照这样的合同条款，全世界2.65万基金持有人在主动管理型公募基金上，投资了1.3万亿美元。

这个世界是怎么做到让9200万美国人参与这样不公平的投资管理协议的？他们付出的管理费相当于未来潜在投资收益的60\%，甚至更多，而自己能够得到的投资收益却没有任何保证。为了解开这个谜局，我坐下来跟投资大师约翰·博格交流，他是先锋公司的创始人。他85岁了，在华尔街干了64年，非常有资格来破解这个难解的金融现象背后的谜团。他的回答是什么呢？

“营销。”

“托尼，其实很简单，大多数人根本没有算账，而且这些费用好多都是隐性的。我们来算算看：如果你20岁一次性投资1万美元，假如每年的投资收益率是7\%，等你到了我这个年龄（80岁），经过60年，你原来投资的1万美元能增值到574464美元。但是，如果你每年要付2.5\%的投资管理费和其他费用，经过同样的60年，你原来投资的1万美元只能增值到140274美元。”

我们来看一下，这笔买卖值不值呢？是你提供所有的投资资本，是你承担所有的投资风险，但你最后只得到了140274美元的投资收益，却付给那些主动管理型基金的基金经理投资管理费439190美元！他们拿走了你潜在收益的77\%？他们凭什么？

“确实如此。”

\textbf{金钱威力法则1：}不懂规则，就不加入比赛。那些基金公司、证券公司、银行等金融机构有组织、有套路地向成千上万的投资者进行营销，以影响投资者的投资决策。我总结为9种投资营销谎言。这些“传统智慧”，经常是故意设计好来蒙骗你的，让你两眼一抹黑，看不清真实的现实。这些传统智慧用在你的金钱投资上，就会蒙住你的双眼。那些你没看到就不知道的事情现在就会，将来可能也会损害你的利益。无知不是什么好事，无知会让你痛苦，无知会让你争斗，无知会让你把财富白白地送给别人，而他什么好处也没有给你。

\subsection{失败的实验}

存在问题的并不只是客户投资成本很高的公募基金。上面讲的例子只是冰山一角，整个金融体系都是故意设计好了的，就是要把你和你的钱分开。

无一例外，我采访的每个专家，从顶尖的基金经理到诺贝尔经济学奖获得者都一致同意，游戏已经变了。和我们现在相比，我们的父辈那个时代的投资理财的复杂性要小得多，危险要少得多。为什么？因为他们有保障生活的退休金，这是由政府保障的终身收入。他们有大额定期储蓄，利率水平不高，但相当合理。我们现在存款利率只有0.22\%，根本赶不上通货膨胀率。有些人还懂点投资，能投资买入一些蓝筹股，股利非常稳定。

船开了，世道变了，退休养老金体制也变了。

美国新的退休养老金体系，在1980年之后才真正开始改革，标志是实施401(k)个人养老投资计划。这其实是美国设计的一场巨大的投资实验，实验人群范围是美国历史上人口最多的一代——婴儿潮世代。这个实验效果怎么样？

“这个新退休养老金体系就是让你自己动手做投资，已经失败了。”特蕾莎·吉拉杜奇说。她是全美国公认的退休养老保障研究的专家，是纽约社会研究新学院的教授，大家都知道她对美国退休养老金体系的批评直言不讳。“失败的原因是，这套退休养老金体系是期望业余投资者获得的收益率能像专业投资者和基金经理一样高。这根本不现实。如果让你自己拔牙，你自己在家里接电线，你想会是什么样的结果？”

是什么变了？我们原来有政府保障的退休养老金，现在被换成了一种新体制，它是全球性的，非常复杂，而且经常是极度危险的，充满了隐性费用。所谓的好处就是能给我们“选择的自由”。在某种程度上，你要辛辛苦苦地工作，养家糊口，锻炼身体、保持身材，要和领导、同事、亲戚、朋友、同学搞好关系，而且你还得成为一个投资专家，做好自己退休养老金的投资管理，你能做到吗？那些设计新的退休养老金体系的政府官员认为，你能够轻松畅游投资迷宫，熟悉了解里面各种各样的金融产品和服务，能够让自己辛辛苦苦赚来的钱避免证券市场源源不断地产生的投资风险。这简直是天方夜谭。正因如此，大部分人都把自己的钱交给“专业投资者”管理，通常交给证券经纪人。证券经纪人是为一家证券公司工作的，在法律上未被要求从最有利于客户的角度出发做事的代理人（我们会在第4个投资营销谎言中做详细分析）。经纪人要想赚钱，靠的是用你的钱去购买那些让他或者他的证券公司赚钱最多的产品。

这里我要先说清楚：本书并不是一本猛烈抨击华尔街——美国金融机构的书。很多大型金融机构已经开创性地设计了一些特别优秀的产品，我们会在本书中仔细研究分析。绝大多数金融服务行业的从业人员非常关心客户，通常情况下，他们做的也是自己认为最好的事情。不幸的是很多人并不明白，这些证券公司、基金公司等金融机构，不管客户是赚是赔，它们自己都要赚钱。这些证券公司尽自己所能为客户做到最好，充分利用它们的知识（培训）和工具（产品）。但是这个体系的设计安排不能让你的证券公司经纪人有无穷的选择，也没有完全的自主权，能够帮你找到最合适的产品，而且要这样做，事实证明，其代价太高了。

你把自己未来潜在收益的很大一部分都作为基金管理费付给基金公司了，或者作为佣金付给证券公司的经纪人了，这非常不合理。但这只是你必须避免的投资陷阱之一。你如果想赢得这场投资游戏比赛，还有很多陷阱也必须避免。不过这里还是有好消息的：

在投资游戏中，你这个业余投资者还是能赢的。

事实上，你不仅能赢，而且投资游戏也会让你非常兴奋。是的，这里有很多困难和挑战，有很多陷阱，你必须尽量避开，但是想一想我们已经走了很远的路。现在只需要按一下按钮，付一点儿费用，你就可以投资全世界任何地方、任何你想要投资的产品。“现在做得比过去好太多了。”詹姆斯·克卢南最近在《华尔街日报》上撰文写道。克卢南是一家非营利机构美国个人投资者大联盟的创办人。“你只需要决定做正确的事情就行了。”

“仅仅在35年前，你还必须花费好几个小时到公共图书馆查阅，或者写信给一家上市公司，只为了能看到它们的财务报告。证券公司的佣金和公募基金的管理费都高得离谱。税收高得简直像偷盗。”贾森·茨威格在《华尔街日报》的一篇专栏文章中写道：“即使股票市场让你紧张，也要想一想你拥有的恩典。”

从科技中受益的不只是高频交易者，科技让整个投资界的效率高了很多，让我们所有的人受益很多。这一点完美地契合了千禧一代的年轻人——任何东西只要慢一点他们就不能接受。“对于我们来说，一切都要快捷方便！”艾米丽大声喊道。她是我的私人助理，她是千禧一代中的一员，说话直来直去。“我们根本无法容忍缓慢和低效率。我们真的希望每一件事情都能够一键搞定。我们买东西都是从亚马逊网上订购的，我们只需要抬一下手指就买好了。我不用下载，就可以在Netflix（视频网站）上播放电影。我可以在网上订购一辆汽车，可以在网上买股票，可以在网上做我的PPT（微软公司的演示文稿软件）课件。今天早上，我拍了一张我的支票的照片，早上6点前支票上的钱就进到了我的银行账户里，我甚至还没有脱掉睡衣呢。”

\subsection{赌场略有优势}

史蒂夫·韦恩，亿万富翁，赌场大亨，他最大的成就就是把拉斯维加斯转变成世界娱乐之都，他也是我的一个最亲密的朋友。史蒂夫·韦恩建造的赌场，大家认为它可以称得上世界上最雄伟的建筑。尽管如此，他赚到巨大财富的秘诀却是一个简单的真理：赌场略有优势，但是这并不能保证他绝对会赢。在一个特定的夜晚，一个连续获胜的赌徒会从史蒂夫的口袋里赢走上百万美元，而且如果他的赌场不能够完全让这些赌客沉迷其中的话，赌客就会转身离开。另一方面，基本上所有的公募基金公司都有优势。它们是超级赌场，它们迷住了你，你哪里也去不了。不管你是赢是输，这些基金公司就像赌场一样，它们的收入都有保证。

\subsection{再燃激情}

2008年之后，美国股市暴跌超过37\%，整个金融世界完全改变了，美国人对金融世界的看法也完全改变了。即使是在五年之后，保德信金融集团的一项调查表明，44\%的美国投资者仍然说他们再也不会往股市里投一分钱了。58\%的人说他们已经对股市丧失了信心。但是那些内行还留在游戏里。为什么？因为他们更懂行了。他们更懂玩这个投资游戏的“正确”方法了，他们也知道今天有很多强有力的工具和策略是自己以前从来没有碰到过的。

现在你可以运用一种金融工具，由全球规模最大的银行发行，并支持以公司资产为你提供100\%的本金保护，而且可以让你分享到市场涨幅的75\%\textasciitilde{}90\%的收益，比如标准普尔500指数。没有上限，这可不是印刷错误，你最高可以分享到市场上涨收益的90\%。但是如果市场下跌了，你还是可以拿回你100\%的本金。这听起来太好了，简直不像是真的。如果确实有这么好的产品存在的话，你应该早就听说过了，对不对？错。为什么？即使你听说过这种产品，你也必须是1\%富豪中的前1\%，你才有机会买到这些金融机构不向大众零售的金融产品，它们不会摆在货架上。这些产品是为那些有足够多财富来参与这种投资的富豪量身定制的。

这只是其中的一个例子。等你成了内行，你很快就会知道新的游戏规则如何用最少的风险获得财富，你就可以了解到更多这样的产品。

\begin{BookQuote}
风险来自你不知道自己正在做什么。

\hfill ——沃伦·巴菲特
\end{BookQuote}

\subsection{少有人走的路}

前面的投资之旅是那种需要你全身心投入来参与的旅程。现在，我们一起去攀登名为财富自由的山峰，这是你个人的珠穆朗玛峰。这条路不会是坦途，需要你提前做准备。你不能一见珠穆朗玛峰就爬，事先不搞清楚前面路上有哪些危险。有些危险是你知道的，有些危险会偷偷接近我们，来个突然袭击，比如一场暴风雪。所以我们开始登山之前，必须充分了解前面路上都有哪些艰难险阻。走错一步，你晚年就要为付下个月的房贷而发愁了；走对了，物质丰富的人生就在前面等着你，让你完全摆脱所有财务压力。我们不能让别人来替我们攀登财务自由的山峰，但是我们也不能一个人做成这件事。我们需要一个向导，他的心里装的永远是我们的最高利益。

\subsection{顶峰}

成功投资的核心理念非常简单：不断储蓄，直到储蓄的资本数量达到一个关键点，之后你的投资收益就能够产生足够多的收入，让你不用工作也能够保持你现在的生活方式，自由自在地生活。最终，你会达到一个引爆点，到那个时候，你的资产规模就达到了一个临界数量级。拥有这么多财富意味着你再也不用去工作了，除非你选择去工作，因为你的投资账户产生的利息和投资升值会给你带来相当多的收入，足以满足你的生活需要。这种财务自由就是我们要攀登的顶峰。一个非常好的消息是，你不断学习进步成了投资内行之后，会有很多新型、独特的解决方案和投资策略，它们会加速你攀登顶峰的速度，甚至可以保护你，避免你滑落下去。但是在更加深入地探索了解这些解决方案前，我们先要清楚地画出我们未来投资之旅的路线图。

你的投资游戏有两个阶段：第一个阶段是积累财富，你要不断地储蓄金钱，追求财富增长；第二个阶段是消耗财富，你不断地提取收入，不断地花钱。上山的旅程代表的是你的财富积累阶段，目标是到达峰顶，也就是到达我们前面所说的临界数量级。我们要尽可能长时间地停留在峰顶，能停多久就停多久。你可以欣赏四周美丽的风光，呼吸新鲜的自由空气，享受自己的伟大成就。这一路上会有很多困难和障碍，如果你不能保持警觉，甚至还对自己撒谎，它们就会阻碍你到达峰顶。为了保证成功的可能性最大，我们要清除这些困难和障碍，后面我们会详述具体办法。

在到达峰顶之后，我们就进入了人生的第二个阶段——消耗财富阶段。这里我们可以享受自己创造的财富，享受我们通过不断地积累才得到的财务自由，从此不再为了钱而工作。你只有在愿意工作的时候才去工作。在这个阶段，我们从山顶一路滑下，只需自由享受自己的人生。跟我们爱的人在一起度过生命宝贵的时光，创建我们想要留给后人的遗产，做与众不同的事。正是在这个阶段，我们会彻底消除婴儿潮世代的最大恐惧：人活着，钱没了。资产管理行业只专注于如何投资赚钱的问题，却很少讨论这个人生的第二阶段，即如何花钱的问题。

“并不是说你到某年某月某日，当你的账户里攒到了足够的钱，你想怎么花就可以怎么花。”杰弗里·布朗博士说。他是伊利诺伊大学的金融学教授，也是美国财政部和世界银行的顾问。“我想，很多人快到退休年龄的时候，会突然醒悟：‘你知道吗？我本来做的工作挺好的，我攒了这么多钱放在那里，但是我不知道我还能活多久，我不知道我的投资收益率会有多高，我不知道未来通货膨胀率会有多高。万一我还活着，钱却花没了，那就惨了。我该怎么办？’”

最近我读到他在《福布斯》杂志专栏中的一篇文章，我马上打电话给布朗博士，看看他是否愿意抽时间坐下来，给各种类型、不同等级的投资者分享一下退休养老金问题的具体解决办法。（在本书的第五部分里，有我采访布朗博士的内容，他会给大家讲解如何创造终身收入，以及如何让这份收入完全免税。）要勾勒出退休养老金问题的解决之道，没有人比布朗博士做得更好了，因为他既是顶尖的学术专家，又是美国总统为社会保障咨询委员会指定的七位顾问之一。

\subsection{挣脱锁链}

戴维·斯文森是我们这个时代最成功的投资者之一，他的忠告是：要想取得不同于传统的成功，就不能接受传统智慧的误导。让我们来打破误导投资大众的9种投资营销谎言，此外，更重要的是去探索发现新的金钱游戏规则，以及那些会让你获得财务自由的投资真理。

我们先从最大的那个投资营销谎言开始……

按照投资百科网站上所说的：“主动管理型的基金经理，依赖于分析、研究、预测，以及他们自己的判断和经验，来做出投资决策——买入、持有、卖出。与之相反的是被动型管理，更通俗的说法是指数投资。
